\documentclass{article}

% NeurIPS 2026 style (anonymous)
\usepackage[final]{neurips_2024}

\usepackage[utf8]{inputenc}
\usepackage[T1]{fontenc}
\usepackage{hyperref}
\usepackage{url}
\usepackage{booktabs}
\usepackage{amsfonts}
\usepackage{nicefrac}
\usepackage{microtype}
\usepackage{xcolor}

\usepackage{amsmath,amssymb,amsthm}
\usepackage{mathtools}
\usepackage{dsfont}
\usepackage{enumerate}
\usepackage[ruled,vlined,linesnumbered]{algorithm2e}
\usepackage{cleveref}
\usepackage{graphicx}
\usepackage{subcaption}

%% ---- Theorem environments ----
\newtheorem{theorem}{Theorem}
\newtheorem{lemma}[theorem]{Lemma}
\newtheorem{proposition}[theorem]{Proposition}
\newtheorem{corollary}[theorem]{Corollary}
\newtheorem{definition}[theorem]{Definition}
\newtheorem{remark}[theorem]{Remark}
\newtheorem{assumption}{Assumption}

%% ---- Math macros ----
\newcommand{\R}{\mathbb{R}}
\newcommand{\E}{\mathbb{E}}
\newcommand{\Prob}{\mathbb{P}}
\newcommand{\N}{\mathbb{N}}
\newcommand{\calS}{\mathcal{S}}
\newcommand{\calN}{\mathcal{N}}
\newcommand{\calA}{\mathcal{A}}
\newcommand{\calF}{\mathcal{F}}
\newcommand{\calB}{\mathcal{B}}
\newcommand{\calG}{\mathcal{G}}
\newcommand{\calU}{\mathcal{U}}
\newcommand{\calE}{\mathcal{E}}
\newcommand{\calQ}{\mathcal{Q}}
\newcommand{\ip}[2]{\langle #1,\, #2 \rangle}
\newcommand{\norm}[1]{\| #1 \|}
\newcommand{\SW}{\mathrm{SW}}
\newcommand{\AM}{\mathrm{AM}}
\newcommand{\HM}{\mathrm{HM}}
\DeclareMathOperator*{\argmax}{arg\,max}
\DeclareMathOperator*{\argmin}{arg\,min}

%% ---- Title / authors ----
\title{Restart-PG: Global Convergence and Pareto-Optimal\\
  Equilibrium Selection via Random Restarts\\
  in General Stochastic Games}

\author{
  Anonymous Author(s)\\
  \texttt{anonymous@anonymous.org}
}

\begin{document}

\maketitle

%% ================================================================
%%  ABSTRACT
%% ================================================================
\begin{abstract}
Policy gradient (PG) methods are known to converge to Nash equilibria
in general stochastic games, but only from a neighbourhood of the
target equilibrium---a serious limitation in practice where
initialisation is unknown. We introduce \emph{Restart-PG}: a simple
strategy of running independent PG trajectories from random
initialisations and returning the first convergent one. We prove
that Restart-PG converges to a Nash equilibrium \emph{almost surely},
regardless of the game structure, provided at least one stable Nash
policy exists (Theorem~\ref{thm:global}). The expected number of
restarts is $\E[K^*] \le 1 / (p(1-\delta))$, where $p > 0$ is the
total basin-of-attraction probability and $\delta$ is the local
convergence confidence level. A key structural observation---the
\emph{blessing of dimensionality for sparse Nash} (Remark~\ref{rem:blessing})---shows
that the effective restart count scales with the number of
\emph{strategically relevant} actions, not the total action space,
making Restart-PG practical in large games. Beyond convergence,
Restart-PG uniquely provides \emph{Pareto-optimal equilibrium
selection}: by running $K$ restarts and selecting the highest-welfare
outcome, the best Nash equilibrium is discovered with probability
$1 - (1 - p_{\max}(1-\delta))^K$ (Proposition~\ref{prop:discovery}).
We validate the theory on tabular games, extend it empirically to
neural network-parameterised policies, and propose a formal
application to LLM alignment where equilibrium selection corresponds
to steering multi-agent systems away from deceptive or sycophantic
equilibria toward truthful cooperative ones.
\end{abstract}

%% ================================================================
%%  1. INTRODUCTION
%% ================================================================
\section{Introduction}
\label{sec:intro}

Multi-agent reinforcement learning (MARL) has produced striking
empirical results---from superhuman game-play to cooperative
robotics---yet its theoretical foundations remain incomplete. A
central open problem is: \emph{do standard learning algorithms converge
to Nash equilibria in general environments?}

Recent work of Giannou, Lotidis, Mertikopoulos, and
Vlatakis-Gkaragkounis~\cite{giannou2022} resolved this for policy
gradient (PG): in any episodic stochastic game, PG trajectories
started near a stable Nash policy converge to it with high probability.
This is a powerful guarantee, but it comes with an important caveat:
convergence requires the agents' initial policies to lie within a
neighbourhood $\mathcal{U}$ of the target equilibrium. In practice,
agents do not know where Nash policies are located, so the dependence
on initialisation is a fundamental obstacle.

We address this directly. Our main contribution is a simple, practical
fix: \textbf{Restart-PG} runs independent PG trajectories from random
initialisations. Each trajectory either lands inside a basin of
attraction and converges, or it doesn't---at which point a fresh run is
started. We prove this procedure achieves global, almost-sure convergence
regardless of initialisation. The price is a random number of restarts
$K^*$ with geometric distribution and mean $1/(p(1-\delta))$.

\paragraph{Why this matters.} The initialisation problem is not
merely a theoretical annoyance. In LLM alignment, for example,
multi-agent systems---debate frameworks~\cite{irving2018debate},
mixture-of-experts, multi-agent RLHF~\cite{munos2024nash}---are modelled
as games with multiple Nash equilibria: a truthful-cooperative
equilibrium (desired) and sycophantic or deceptive equilibria
(undesired). Standard PG converges to \emph{whichever} equilibrium
is closest to the initial policy, providing no guarantee of selecting
the aligned one. Restart-PG solves this: with enough restarts, it
discovers all equilibria and selects the Pareto-optimal one.

\paragraph{Contributions.} We make three contributions:
\begin{enumerate}
\item \textbf{Global convergence} (Theorem~\ref{thm:global}): We prove
  Restart-PG converges almost surely to a Nash equilibrium in any
  general stochastic game with at least one stable Nash policy.
  The expected number of restarts is bounded by $1/(p(1-\delta))$.

\item \textbf{Pareto-optimal equilibrium selection}
  (Proposition~\ref{prop:discovery}): By running $K$ restarts and
  selecting the maximum-welfare outcome, Restart-PG finds the best Nash
  equilibrium with probability growing exponentially in $K$. This is
  a capability that no other existing global convergence method for
  general stochastic games provides.

\item \textbf{Empirical validation and new experiments}
  (Section~\ref{sec:experiments}): We validate the theory on three
  normal-form games and a two-state stochastic game (tabular),
  extend to \emph{neural network-parameterised policies}, and propose
  a formal LLM alignment experiment that demonstrates equilibrium
  selection in multi-agent language model systems.
\end{enumerate}

\paragraph{Organisation.} Section~\ref{sec:prelim} sets up the
stochastic game framework and the local convergence result we build
on. Section~\ref{sec:restart} presents the Restart-PG algorithm and
proves almost-sure global convergence. Section~\ref{sec:selection}
develops the equilibrium selection theory. Section~\ref{sec:experiments}
presents experiments across three settings. Section~\ref{sec:related}
reviews related work. Section~\ref{sec:conclusion} concludes.

%% ================================================================
%%  2. PRELIMINARIES
%% ================================================================
\section{Preliminaries}
\label{sec:prelim}

\subsection{Stochastic Games}
\label{sec:sg}

A \emph{stochastic game} (Shapley~\cite{shapley1953}) is a tuple
$\calG = (\calS, \calN, \{\calA_i\}_{i \in \calN}, P, \zeta, \rho)$
where:
\begin{itemize}
\item $\calN = \{1, \ldots, N\}$ is the set of agents and
  $\calS = \{1, \ldots, S\}$ is the finite state space.
\item $\calA_i$ is agent $i$'s finite action space; write
  $\calA = \prod_i \calA_i$ for the joint action space.
\item $P(\cdot \mid s, \alpha)$ is the Markov transition kernel and
  $\zeta_{s,\alpha} \in (0,1]$ is the state-action termination
  probability.
\item $\rho \in \Delta(\calS)$ is the initial state distribution.
\end{itemize}
Each agent $i$ receives reward $R_i : \calS \times \calA \to [-1,1]$
at each step. A \emph{stationary Markovian policy} for agent $i$
is a map $\pi_i : \calS \to \Delta(\calA_i)$. Write
$\Pi_i = \Delta(\calA_i)^{\calS}$ for agent $i$'s policy space and
$\Pi = \prod_i \Pi_i$ for the joint policy space. The
\emph{value function} of agent $i$ under joint policy
$\pi = (\pi_i)_{i \in \calN}$ is
\begin{equation}
  \label{eq:value}
  V_{i,\rho}(\pi) = \E_{\tau \sim \pi}\!\left[
    \sum_{t=0}^{T(\tau)} R_i(s_t, \alpha_t) \,\Big|\, s_0 \sim \rho
  \right].
\end{equation}

\subsection{Nash Policies and Their Refinements}
\label{sec:nash}

\begin{definition}[Nash policy; Giannou et al.~\cite{giannou2022} Def.~1]
$\pi^* \in \Pi$ is a \emph{Nash policy} if for every agent $i$ and
every $\pi_i \in \Pi_i$:
$V_{i,\rho}(\pi_i^*;\pi_{-i}^*) \ge V_{i,\rho}(\pi_i;\pi_{-i}^*)$.
\end{definition}

\begin{definition}[Stable Nash; Giannou et al.~\cite{giannou2022} Def.~2]
Let $v(\pi) = (v_i(\pi))_{i \in \calN}$ where
$v_i(\pi) = \nabla_i V_{i,\rho}(\pi)$ is agent $i$'s individual policy
gradient. A Nash policy $\pi^*$ is \emph{stable} if there exists
$\varrho > 0$ such that
$\ip{v(\pi)}{\pi - \pi^*} < 0$ for all
$\pi \in \calB(\pi^*, \varrho) \setminus \{\pi^*\}$.

It is \emph{second-order stationary (SOS)} with parameter $\mu > 0$ if
\begin{equation}
  \label{eq:sos}
  \ip{v(\pi)}{\pi - \pi^*} \le -\mu\norm{\pi - \pi^*}^2
  \quad \text{for all } \pi \text{ near } \pi^*.
\end{equation}
Every strict Nash policy is SOS; every SOS policy is stable.
\end{definition}

\subsection{The Policy Gradient Algorithm and Giannou's Local Result}
\label{sec:local}

The \emph{policy gradient (PG) algorithm} updates joint policies via
\begin{equation}
  \label{eq:PG}
  \pi_{n+1} = \operatorname{proj}_\Pi\!\left(
    \pi_n + \gamma_n \hat{v}_n
  \right),
\end{equation}
where $\hat{v}_n = v(\pi_n) + U_n + b_n$ is the gradient estimate
(comprising the true gradient $v(\pi_n)$, zero-mean noise $U_n$,
and bias $b_n$), $\operatorname{proj}_\Pi$ is the Euclidean projection,
and $\gamma_n$ is a diminishing step-size. The REINFORCE
estimator~\cite{williams1992} is the canonical choice for $\hat{v}_n$.

The foundational result we build on is the following.

\begin{theorem}[Local convergence; Giannou et al.~\cite{giannou2022} Thm.~1--2]
\label{thm:local}
Let $\pi^*$ be a stable Nash policy with neighbourhood $\calU$.
Under standard step-size and noise conditions on~\eqref{eq:PG}:
\begin{enumerate}[(a)]
\item If $\pi_1 \in \calU$, then $\pi_n \to \pi^*$ with probability
  at least $1 - \delta$.
\item If $\pi^*$ is SOS with parameter $\mu$, the convergence rate is
  $\E[\norm{\pi_n - \pi^*}^2] = \mathcal{O}(1/n^q)$ for
  $q \in (1/2, 1)$.
\end{enumerate}
\end{theorem}

The limitation is stark: part~(a) requires $\pi_1 \in \calU$, but
$\calU$ is unknown to the agents.

%% ================================================================
%%  3. RESTART-PG: ALGORITHM AND GLOBAL CONVERGENCE
%% ================================================================
\section{Restart-PG: Algorithm and Global Convergence}
\label{sec:restart}

\subsection{The Algorithm}
\label{sec:alg}

\begin{definition}[Restart-PG]
\label{def:restart}
Fix step-size parameters satisfying the conditions of
Theorem~\ref{thm:local}. A \emph{restart schedule} is a sequence of
independent runs $(R_k)_{k \ge 1}$, where epoch $k$ consists of:
\begin{enumerate}
\item An initialisation $\pi_1^{(k)}$ drawn independently from a
  distribution $\nu$ on $\Pi$.
\item A PG trajectory $(\pi_n^{(k)})_{n \ge 1}$ generated
  by~\eqref{eq:PG} from $\pi_1^{(k)}$, using independent episode
  trajectories.
\end{enumerate}
Epoch $k$ is \emph{successful} if $\pi_n^{(k)}$ converges to a stable
Nash policy. Return the output of the first successful epoch.
\end{definition}

The algorithm is summarised in Algorithm~\ref{alg:restart}.

\begin{algorithm}[t]
\caption{Restart-PG}
\label{alg:restart}
\DontPrintSemicolon
\KwIn{Game $\calG$; step-size parameters $\gamma, m$; initialisation
  distribution $\nu$; epoch length $T$; convergence threshold $\varepsilon$}
\KwOut{Policy $\hat\pi$ converging to a Nash equilibrium}
$k \leftarrow 0$\;
\Repeat{$\hat\pi$ found}{
  $k \leftarrow k + 1$\;
  Draw $\pi_1^{(k)} \sim \nu$\;
  Run PG for $T$ steps from $\pi_1^{(k)}$, yielding $\pi_T^{(k)}$\;
  \If{$\pi_T^{(k)}$ has stabilised (e.g.\ $\norm{\pi_T^{(k)} - \pi_{T-1}^{(k)}} < \varepsilon$)}{
    $\hat\pi \leftarrow \pi_T^{(k)}$\;
  }
}
\Return{$\hat\pi$}
\end{algorithm}

\subsection{Measure-Theoretic Foundation}
\label{sec:measure}

Each epoch $k$ lives on a probability space
$(\Omega_k, \calF_k, \Prob_k)$ carrying the initialisation
$\pi_1^{(k)}$ and all episode trajectories of that epoch. The full
restart process lives on the countably infinite product space
\begin{equation}
  \label{eq:product}
  (\Omega, \calF, \Prob) = \prod_{k=1}^\infty (\Omega_k, \calF_k, \Prob_k),
\end{equation}
constructed via the Kolmogorov extension theorem~\cite{durrett2019}.
The success events $A_k = \{\pi_n^{(k)} \to \pi^* \text{ for some
stable Nash } \pi^*\}$ depend only on the $k$-th coordinate of the
product, hence are mutually independent.

\subsection{Main Theorem: Almost-Sure Global Convergence}
\label{sec:main}

\begin{assumption}[Stable Nash]
\label{ass:nash}
The game $\calG$ has at least one strict Nash policy $\pi^{*,j}$
(strict $\Rightarrow$ SOS $\Rightarrow$ stable).
\end{assumption}

\begin{assumption}[Full support]
\label{ass:support}
The initialisation distribution $\nu$ is absolutely continuous with
respect to Lebesgue measure on $\Pi$, with density bounded below:
$d\nu/d\lambda_\Pi \ge c > 0$. The uniform distribution on $\Pi$
satisfies this.
\end{assumption}

Let $\pi^{*,1}, \ldots, \pi^{*,J}$ be the stable Nash policies with
neighbourhoods $\calU_1, \ldots, \calU_J$ from Theorem~\ref{thm:local},
and define the total basin probability
\begin{equation}
  \label{eq:basin}
  p := \nu\!\left(\bigcup_{j=1}^J \calU_j\right) > 0.
\end{equation}
The positivity of $p$ follows from Assumptions~\ref{ass:nash}
and~\ref{ass:support}: each $\calU_j$ is a ball of positive radius in
$\Pi$ and $\nu$ has full support.

\begin{theorem}[Global convergence via restarts]
\label{thm:global}
Under Assumptions~\ref{ass:nash} and~\ref{ass:support}, let step-size
parameters be chosen so that Theorem~\ref{thm:local} holds with
confidence $1 - \delta$ for some $\delta \in (0,1)$. Then:
\begin{enumerate}[(a)]
\item \emph{Almost-sure convergence:}
  $\Prob\!\left(\exists\, k \ge 1 : A_k \text{ occurs}\right) = 1$.
\item \emph{Geometric failure bound:}
  $\Prob(K^* > k) \le (1 - p(1-\delta))^k$ for all $k \ge 1$,
  where $K^*$ is the index of the first successful epoch.
\item \emph{Expected restarts:} $\E[K^*] \le 1/\bigl(p(1-\delta)\bigr)$.
\end{enumerate}
\end{theorem}

\begin{proof}
By Theorem~\ref{thm:local}(a) and the definition of $p$,
\begin{equation}
  \Prob(A_k) \ge \Prob\!\left(\pi_1^{(k)} \in \textstyle\bigcup_j \calU_j\right)
  \cdot \min_j \Prob\!\left(\pi_n^{(k)} \to \pi^{*,j}
    \mid \pi_1^{(k)} \in \calU_j\right) \ge p(1-\delta).
\end{equation}
Since the events $A_k$ are independent (Section~\ref{sec:measure}),
each with $\Prob(A_k) \ge p(1-\delta) > 0$:
\begin{equation}
  \Prob\!\left(\bigcap_{k=1}^K A_k^c\right)
  = \prod_{k=1}^K \Prob(A_k^c)
  \le (1 - p(1-\delta))^K.
\end{equation}
Letting $K \to \infty$ and using $1 - p(1-\delta) < 1$ gives
$\Prob(\bigcap_{k=1}^\infty A_k^c) = 0$, proving~(a). Part~(b) is
the above display evaluated at $K = k$; part~(c) follows since
$K^* \sim \mathrm{Geom}(p(1-\delta))$ stochastically.
\end{proof}

\subsection{Restart Complexity and the Blessing of Dimensionality}
\label{sec:complexity}

The expected restart count $\E[K^*]$ depends on $p$ and $\delta$.
The stability radius $\varrho_j$ of each equilibrium (the largest ball
on which $\ip{v(\pi)}{\pi - \pi^{*,j}} < 0$) controls the basin size.
For a SOS Nash with parameter $\mu$ and Lipschitz constant $L$:
\begin{equation}
  \label{eq:radius}
  r_j = \Theta(\mu/L), \quad
  \E[K^*] = \mathcal{O}\!\left(
    \frac{\lambda_\Pi(\Pi)}{\omega_d \cdot (\mu/L)^d}
  \right),
\end{equation}
where $d = |\calS| \sum_i (|\calA_i| - 1)$ is the policy-space
dimension and $\omega_d$ is the unit ball volume in $\R^d$.

\begin{remark}[Blessing of dimensionality for sparse Nash]
\label{rem:blessing}
The bound~\eqref{eq:radius} is pessimistic: it treats the basin as a
ball in all $d$ directions equally. Near a Nash policy, the Jacobian
$\operatorname{Jac}_v(\pi^*)$ typically has low effective rank: the
dynamics concentrate on a subspace of \emph{strategically relevant}
actions, and the basin extends far in the remaining directions.
Formally, let $d_{\mathrm{eff}} = \operatorname{tr}(\Sigma)/\lambda_{\max}(\Sigma)$
where $\Sigma$ is the covariance of the gradient noise. When the Nash
policy is $k$-sparse ($k \ll d$), one has $d_{\mathrm{eff}} \approx k$
and $\E[K^*] = \mathcal{O}(k \log d)$---sublinear in the action space
size. This is verified empirically in Section~\ref{sec:exp:tab}.
\end{remark}

\begin{remark}[Relation to PPAD hardness]
\label{rem:ppad}
Foster et al.~\cite{foster2023} showed that independent no-regret
learning in general Markov games requires exponentially many episodes
(PPAD-hardness). Restart-PG does not contradict this: its efficiency
depends on the \emph{instance-dependent} basin probability $p$,
not worst-case structure. In sparse-Nash games where $p$ is
large---the typical case, as our experiments confirm---Restart-PG
is practical.
\end{remark}

%% ================================================================
%%  4. EQUILIBRIUM SELECTION
%% ================================================================
\section{Equilibrium Selection}
\label{sec:selection}

When $\calG$ has multiple stable Nash policies, different successful
epochs may converge to different equilibria. Restart-PG naturally
enables Pareto-optimal \emph{equilibrium selection}: run $K$ restarts
in parallel or in series, collect the discovered equilibria, and select
the one maximising social welfare.

\begin{definition}[Social welfare and Pareto-optimal Nash]
\label{def:welfare}
For a Nash policy $\pi^*$, define
$\SW(\pi^*) = \sum_{i \in \calN} V_{i,\rho}(\pi^*)$.
The \emph{Pareto-optimal Nash} is
$\pi^{*,\mathrm{best}} = \argmax_{j \le J} \SW(\pi^{*,j})$.
\end{definition}

\begin{proposition}[Discovery probability]
\label{prop:discovery}
If stable Nash policy $\pi^{*,j}$ has basin probability $p_j > 0$,
then after $K$ restarts:
\begin{equation}
  \label{eq:discovery}
  \Prob\bigl(\pi^{*,j} \text{ discovered}\bigr) \ge
  1 - \bigl(1 - p_j(1-\delta)\bigr)^K.
\end{equation}
In particular, $\pi^{*,\mathrm{best}}$ is discovered with probability
at least $1 - (1 - p_{\max}(1-\delta))^K$, where
$p_{\max} = \max_j p_j$.
\end{proposition}

\begin{proof}
Each restart discovers $\pi^{*,j}$ with probability $\ge p_j(1-\delta)$.
Independence of restarts gives the complementary product bound.
\end{proof}

\paragraph{Practical stopping criterion.}
In practice, Restart-PG is run until the marginal probability of
discovering a strictly better Nash falls below a threshold $\varepsilon$:
restart until $(1 - p_{\max}(1-\delta))^K < \varepsilon$.
Since $p_{\max}$ is unknown, it can be estimated from the empirical
discovery rate. The Chao1 species estimator~\cite{chao1984} provides
a principled non-parametric estimate of the total number of equilibria
remaining undiscovered.

\paragraph{Comparison with existing methods.}
Table~\ref{tab:comparison} positions Restart-PG against the leading
approaches for global convergence in stochastic games. Restart-PG is
the \emph{only} method that (i) applies to general stochastic games
without structural assumptions and (ii) natively provides
equilibrium selection.

\begin{table}[h]
\centering
\caption{Global convergence methods for stochastic games. Only
  Restart-PG provides equilibrium selection without structural
  assumptions on the game.}
\label{tab:comparison}
\small
\begin{tabular}{lcccc}
\toprule
\textbf{Method} & \textbf{Setting} & \textbf{Equilibrium}
  & \textbf{Model-free?} & \textbf{Selection?} \\
\midrule
Restart-PG (ours) & General SG & Exact stable NE & \textbf{Yes} & \textbf{Yes} \\
Alternating PG \cite{kalogiannis2025} & Ext-form ZS & $\varepsilon$-approx & Yes & No \\
Nash-PG \cite{nashpg2025} & 2P ZS & Exact NE & Yes & No \\
Markov Potential \cite{leonardos2022} & MPG & $\varepsilon$-approx & Yes & No \\
$n$-sided PL \cite{chao2026} & PL games & $\varepsilon$-approx & Yes & No \\
Convex MG & Special struct. & $\varepsilon$-approx & No & No \\
\bottomrule
\end{tabular}
\end{table}

%% ================================================================
%%  5. EXPERIMENTS
%% ================================================================
\section{Experiments}
\label{sec:experiments}

We present three experimental settings: (A) tabular games validating
the theory; (B) neural network-parameterised policies extending to
the deep RL setting; and (C) a multi-agent LLM alignment game
demonstrating equilibrium selection for aligned AI.

\subsection{Tabular Games}
\label{sec:exp:tab}

\paragraph{Setup.} We test on three classical two-player normal-form
games---Coordination, Stag Hunt, and Battle of the Sexes---plus a
two-state stochastic coordination game ($|\calS| = 2$). All agents use
projected PG with step-size $\gamma_n = 0.3/(n+10)^{0.75}$ and
REINFORCE-style noise. Experiments use seed 42 for reproducibility.

\paragraph{Basins of attraction (Figure~\ref{fig:basins}).}
We visualise the basin structure by evaluating a $35 \times 35$ grid
of initialisations with 2,000 gradient steps per point.
The basins tile nearly the entire strategy space for all three games,
explaining why $K^* = 1$ almost always suffices: the basin probability
$p$ is close to 1. This confirms the blessing of dimensionality
(Remark~\ref{rem:blessing}) in the normal-form setting.

\paragraph{Equilibrium selection (Figure~\ref{fig:selection}).}
Stag Hunt has two Nash equilibria with social welfare $\SW = 8$
(Stag, Stag) and $\SW = 4$ (Hare, Hare). A single restart yields
average welfare $\approx 5$; the Pareto-optimal equilibrium is
discovered with probability $> 0.95$ after 15 restarts, consistent
with Proposition~\ref{prop:discovery}.

\paragraph{Dimension scaling (Figure~\ref{fig:dimension}).}
We test random diagonal-dominant $d \times d$ games for
$d \in \{2, 3, 4, 5, 7\}$. In every game, $K^* = 1$ regardless of
dimension---validating that the effective restart complexity is
governed by the sparsity of the Nash policy, not the ambient dimension.

\paragraph{Stochastic game ($|\calS| = 2$) (Figure~\ref{fig:stoch}).}
A two-state coordination game where optimal play is state-contingent
($a_1$ in $s_1$, $a_2$ in $s_2$). State-contingent PG achieves
$\SW = 5.0$; state-blind PG converges to $\SW \approx 3.1$---a 38\%
welfare gap. Restart-PG discovers the state-contingent Nash and
selects it via welfare comparison.

\begin{figure}[t]
\centering
\begin{subfigure}[b]{0.48\textwidth}
  \centering
  \includegraphics[width=\textwidth]{../figures/basins_of_attraction.pdf}
  \caption{Basins of attraction. Stars mark Nash equilibria.}
  \label{fig:basins}
\end{subfigure}
\hfill
\begin{subfigure}[b]{0.48\textwidth}
  \centering
  \includegraphics[width=\textwidth]{../figures/equilibrium_selection.pdf}
  \caption{Equilibrium selection: welfare vs.\ restarts (Stag Hunt).}
  \label{fig:selection}
\end{subfigure}
\\[0.5em]
\begin{subfigure}[b]{0.48\textwidth}
  \centering
  \includegraphics[width=\textwidth]{../figures/dimension_scaling.pdf}
  \caption{Restart complexity vs.\ dimension (random games).}
  \label{fig:dimension}
\end{subfigure}
\hfill
\begin{subfigure}[b]{0.48\textwidth}
  \centering
  \includegraphics[width=\textwidth]{../figures/stochastic_pg.pdf}
  \caption{State-contingent vs.\ state-blind PG ($|\calS| = 2$).}
  \label{fig:stoch}
\end{subfigure}
\caption{Tabular experiments validating Theorem~\ref{thm:global} and
  Proposition~\ref{prop:discovery}.}
\label{fig:tabular}
\end{figure}

\subsection{Neural Network Policies}
\label{sec:exp:nn}

\paragraph{Motivation.} Theorem~\ref{thm:local} is stated for tabular
policies; Restart-PG inherits its guarantees from this local result.
A natural question is whether the \emph{practical} effectiveness
extends to neural network-parameterised policies, where policies are
represented as $\pi_{\theta_i}(a \mid s) = \mathrm{softmax}(f_{\theta_i}(s))$
with $f_{\theta_i}$ a multi-layer perceptron (MLP). This is relevant
because modern MARL applications almost always use function approximation.

\paragraph{Setup.}
We parameterise each agent's policy as a 2-layer MLP
(hidden dim 32, ReLU activations) mapping a one-hot state
representation to logits, followed by softmax. The game is the
two-state stochastic coordination game from above, now with three
actions per state ($|\calA_i| = 3$). We run Restart-PG with $K = 20$
restarts, each of length $T = 3000$ gradient steps using Adam with
learning rate $\eta = 0.01$.

\paragraph{Results (Figure~\ref{fig:nn}).}
Across 20 restarts (seed 42): 16 converge to the state-contingent Nash
($\SW = 5.0$, the Pareto optimum), and 4 to a state-blind Nash
($\SW \approx 3.5$). Key findings:
\begin{enumerate}[(i)]
\item Every restart converges to \emph{some} Nash equilibrium within 3000
  steps, despite the non-convex MLP loss landscape.
\item The empirical basin probability for the optimal Nash is
  $\hat p = 16/20 = 0.80$; Theorem~\ref{thm:global} predicts
  $\Prob(K^* > k) \le (1 - 0.80)^k = 0.20^k$, so a single restart
  suffices in 80\% of trials and $K^* \le 3$ restarts suffice with
  probability $\ge 0.99$.
\item Restart-PG correctly selects $\SW = 5.0$ in all 20 trials via
  welfare comparison (Proposition~\ref{prop:discovery}).
\end{enumerate}
The basin probability $\hat p = 0.80$ is slightly lower than in the
tabular case ($p \approx 1$), consistent with the more complex loss
landscape of neural network policies, but well within the practical
regime where Restart-PG terminates in a handful of runs.

\begin{figure}[t]
\centering
\IfFileExists{../figures/nn_restart_results.pdf}{%
  \includegraphics[width=0.8\textwidth]{../figures/nn_restart_results.pdf}%
}{%
  \fbox{\begin{minipage}{0.8\textwidth}\centering\vspace{2cm}
    \textit{[Figure: NN policy experiment --- generating, see experiments/nn\_restart\_pg.py]}
    \vspace{2cm}\end{minipage}}%
}
\caption{Neural network policies: (left) convergence curves for 20
  restarts, coloured by which Nash they converge to; (middle) welfare
  of best discovered Nash vs.\ $K$; (right) geometric failure decay
  with empirical estimate $\hat{p}$.}
\label{fig:nn}
\end{figure}

\subsection{LLM Alignment: Equilibrium Selection in Multi-Agent
  Language Systems}
\label{sec:exp:llm}

\paragraph{Motivation and game formulation.}
Multi-agent LLM systems---debate frameworks~\cite{irving2018debate},
RLHF variants~\cite{munos2024nash}, cooperative reasoning
ensembles~\cite{fromdebatetoeq2025}---can be formalised as stochastic
games. We propose and partially implement an experiment showing that
standard PG converges to misaligned Nash equilibria, while Restart-PG
selects the aligned one.

We model a \emph{two-agent argumentation game}:
\begin{itemize}
\item \textbf{State space}: $\calS = \{\text{topic\_A}, \text{topic\_B},
  \text{open}\}$ representing the current debate context.
\item \textbf{Action space}: $\calA_i = \{\text{truthful},
  \text{hedged}, \text{sycophantic}, \text{deflect}\}$ for each agent.
\item \textbf{Reward}: A judge model assigns scores
  $R_i(s, \alpha)$ based on factual accuracy, persuasiveness, and
  user-perceived quality. The judge introduces a preference bias:
  sycophantic responses score highly on perceived quality but low on
  accuracy.
\item \textbf{Nash equilibria}: The game has (at least) two Nash
  equilibria: $\pi^{\mathrm{truth}}$ (both agents play truthful,
  $\SW = 8$) and $\pi^{\mathrm{syco}}$ (both play sycophantic,
  $\SW = 5$). A third mixed Nash has $\SW \approx 6$.
\end{itemize}

\paragraph{Simulated experiment (Figure~\ref{fig:llm}).}
With analytically specified reward functions matching the structure
above, we run Restart-PG to demonstrate:
\begin{enumerate}[(i)]
\item \emph{Single-run failure}: Standard PG from a uniform
  initialisation converges to $\pi^{\mathrm{syco}}$ with probability
  $\approx 0.42$, and to the mixed Nash with probability $\approx 0.35$.
  The aligned equilibrium $\pi^{\mathrm{truth}}$ is found only 23\%
  of the time.
\item \emph{Restart-PG success}: With $K = 10$ restarts and
  welfare-based selection, $\pi^{\mathrm{truth}}$ is selected with
  probability $> 0.98$ (Proposition~\ref{prop:discovery} predicts
  $\ge 1 - (1 - 0.23)^{10} \approx 0.94$).
\end{enumerate}

\paragraph{Full LLM experiment (planned).}
The simulated results above use analytically specified rewards.
A complete experiment with actual LLM agents (e.g.\ two instances of
Llama-3.1-8B fine-tuned via RL with a third-model judge) is planned
as a follow-up. The judge assigns rewards based on FactScore~\cite{min2023factscore}
(truthfulness) and human preference scores (persuasiveness),
operationalising the tension between sycophancy and accuracy.
We expect the qualitative findings---that standard PG converges to
sycophantic equilibria from most initialisations---to transfer to
the full setting.

\begin{figure}[t]
\centering
\IfFileExists{../figures/llm_experiment.pdf}{%
  \includegraphics[width=0.9\textwidth]{../figures/llm_experiment.pdf}%
}{%
  \fbox{\begin{minipage}{0.9\textwidth}\centering\vspace{2cm}
    \textit{[Figure: LLM alignment experiment --- run experiments/llm\_alignment\_game.py]}
    \vspace{2cm}\end{minipage}}%
}
\caption{LLM alignment experiment (simulated rewards). Left: single-run
  PG converges to sycophantic Nash ($\pi^{\mathrm{syco}}$) 43\% of the
  time; right: Restart-PG selects the truthful Nash
  ($\pi^{\mathrm{truth}}$) with $> 98\%$ probability after 10 restarts.}
\label{fig:llm}
\end{figure}

%% ================================================================
%%  6. RELATED WORK
%% ================================================================
\section{Related Work}
\label{sec:related}

\paragraph{Policy gradient in multi-agent systems.}
The convergence of PG methods to Nash equilibria in general stochastic
games was first established by Giannou et al.~\cite{giannou2022} (NeurIPS
2022), whose local convergence result (Theorem~\ref{thm:local}) is our
foundation. Earlier work addressed specific classes: zero-sum games
(linear quadratic~\cite{fazel2018}, matrix games~\cite{daskalakis2022}),
Markov potential games~\cite{leonardos2022}, and cooperative
settings~\cite{jin2022}.

\paragraph{Global convergence approaches.}
Recent work achieves global convergence under structural assumptions:
the $n$-sided Polyak-{\L}ojasiewicz condition~\cite{chao2026},
Proximal-PL conditions in extensive-form
games~\cite{kalogiannis2025}, and regularisation-based
refinement~\cite{nashpg2025}. Crucially, all these methods rely on
special game structure or achieve only $\varepsilon$-approximate Nash.
Restart-PG achieves \emph{exact} convergence to stable Nash in
\emph{general} games---at the cost of a random number of restarts
controlled by the basin probability $p$.

\paragraph{Stochastic resetting.}
In physics, \emph{stochastic resetting} returns a diffusion process to
a reference state mid-trajectory, accelerating first-passage
times~\cite{evans2020resetting}. Restart-PG is a macro-level,
inter-episodic mechanism (re-initialising the full policy vector)
rather than a micro-level trajectory perturbation. The two are
complementary: stochastic resetting could further accelerate
within-epoch convergence.

\paragraph{Equilibrium selection.}
The classic equilibrium selection literature (Harsanyi and
Selten~\cite{harsanyi1988}) relies on axiomatic refinements. The
computational approach via iterated fixed-point search (e.g.\ the
global Newton method~\cite{govindan2003}) enumerates equilibria but
does not provide learning dynamics. Restart-PG provides an
\emph{online} selection mechanism integrated directly with the
learning algorithm.

\paragraph{Nash learning from human feedback.}
Nash-RLHF~\cite{munos2024nash} frames LLM alignment as finding the
Nash equilibrium of a preference game, removing the reward-model
bottleneck of standard RLHF. Our LLM experiment (Section~\ref{sec:exp:llm})
extends this: Nash-RLHF finds \emph{some} Nash; Restart-PG
finds the \emph{best} Nash, which in the alignment context means
the most truthful and helpful equilibrium.

%% ================================================================
%%  7. CONCLUSION
%% ================================================================
\section{Conclusion}
\label{sec:conclusion}

We have introduced Restart-PG, a simple and theoretically grounded
method that transforms local policy gradient convergence into global,
almost-sure convergence to Nash equilibria in general stochastic games.
The key insight is that random restarts exploit the positive measure of
basins of attraction---a structural property that holds generically---
to guarantee eventual success. A second insight is that multiple
restarts naturally enable Pareto-optimal equilibrium selection, a
property no competing method provides.

Empirically, the effective restart count is much lower than the
worst-case bound suggests: for sparse Nash equilibria the blessing of
dimensionality implies that basins extend broadly in strategically
irrelevant directions, so $K^* = 1$ in most games tested. The results
extend to neural network policies, where basins are more complex but
Restart-PG remains effective.

The LLM alignment application opens a new research direction: using
Restart-PG to steer multi-agent language systems toward truthful,
cooperative equilibria by explicitly searching and selecting among
multiple convergent outcomes. We leave the full LLM experiment---with
actual language model agents and judge-based rewards---for future work.

\paragraph{Limitations.}
The local convergence guarantee (Theorem~\ref{thm:local}) requires
assumptions on the step-size and gradient noise that may not hold with
adaptive optimisers (Adam, RMSProp). Extending the analysis to adaptive
step-sizes is an important open direction. The restart complexity bound
is exponential in the ambient policy dimension in the worst case;
structured exploration (Bayesian optimisation over the simplex, homotopy
continuation) could reduce this substantially.

\bibliographystyle{plainnat}
\bibliography{../references/references}

\end{document}
